% =================================================================
% LaTeX header - Articulo Universitario UTP
% Aplica: colores UTP, callouts EV-XX como cajas, tablas compactas,
%         saltos de pagina por seccion, fonts con fallback
% =================================================================

% --- Colores institucionales UTP ---
\usepackage{xcolor}
\definecolor{utpred}{HTML}{A6192E}
\definecolor{utpredlight}{HTML}{F0E6E8}
\definecolor{calloutbg}{HTML}{FFF8E1}
\definecolor{calloutborder}{HTML}{F9A825}
\definecolor{callouttext}{HTML}{5D4037}

% --- Titulos coloreados en rojo UTP ---
\usepackage{titlesec}
\titleformat{\section}
  {\Large\bfseries\color{utpred}}
  {}{0pt}
  {}
  [\vspace{-0.4em}\color{utpred}\rule{\linewidth}{1pt}]
\titleformat{\subsection}
  {\large\bfseries\color{utpred!85!black}}
  {}{0pt}{}
\titleformat{\subsubsection}
  {\normalsize\bfseries\color{utpred!70!black}}
  {}{0pt}{}

% --- Apartados romanos (I, II, III, IV...) inician en pagina nueva ---
% Los apartados estan como H2 (## en markdown) -> \subsection en LaTeX.
% Tras eliminar el titulo duplicado del cuerpo, la secuencia de H2 es:
%   #1 ABSTRACT, #2 KEYWORDS, #3 I. INTRODUCCION, #4 II. ..., ...
% ABSTRACT + KEYWORDS van juntos despues de la portada; desde I. en adelante
% cada apartado romano abre pagina nueva.
\newcounter{h2count}
\setcounter{h2count}{0}
\let\oldsubsection\subsection
\renewcommand{\subsection}{%
  \stepcounter{h2count}%
  \ifnum\value{h2count}>2\relax%
    \clearpage%
  \fi
  \oldsubsection%
}

% --- Tablas mas compactas (ajuste de fuente y separacion) ---
\usepackage{etoolbox}
\usepackage{array}
\usepackage{ragged2e}     % \RaggedRight para que las celdas hagan wrap
\AtBeginEnvironment{longtable}{\footnotesize\RaggedRight}
\renewcommand{\arraystretch}{1.2}
\setlength{\tabcolsep}{4pt}    % menos padding entre columnas (default 6pt)
\setlength{\LTleft}{0pt}        % alinear longtable a la izquierda del margen
\setlength{\LTright}{0pt}

% --- Forzar a longtable a tipar las columnas como parrafo (no left fijo)
% para que el texto haga wrap automatico al ancho disponible.
% Truco: redefinir el codigo de columna 'l' como 'p{...}' dentro de longtable.
\AtBeginEnvironment{longtable}{%
  \renewcommand{\arraystretch}{1.2}%
}

% --- Bordes mas claros para tablas tipo booktabs ---
\usepackage{booktabs}
\usepackage{colortbl}    % requerido para \arrayrulecolor
\setlength{\heavyrulewidth}{1.5pt}
\setlength{\lightrulewidth}{0.6pt}
\arrayrulecolor{utpred!50!black}

% --- Callouts EV-XX: convertir blockquotes en cajas resaltadas ---
\usepackage[most]{tcolorbox}
\renewenvironment{quote}{%
  \begin{tcolorbox}[
    colback=calloutbg,
    colframe=calloutborder,
    coltext=callouttext,
    leftrule=4pt,
    toprule=0pt,
    rightrule=0pt,
    bottomrule=0pt,
    arc=2pt,
    top=6pt,
    bottom=6pt,
    left=10pt,
    right=10pt,
    breakable
  ]%
}{%
  \end{tcolorbox}%
}

% --- Hipervinculos en rojo UTP ---
\usepackage{hyperref}
\hypersetup{
  colorlinks=true,
  linkcolor=utpred,
  urlcolor=utpred,
  citecolor=utpred,
  pdftitle={Articulo Academico UTP - Redes y Comunicacion de Datos I},
  pdfauthor={Vega, Benites, Quispe, Tasayco}
}

% --- Numero de pagina en pie ---
\usepackage{fancyhdr}
\pagestyle{fancy}
\fancyhf{}
\fancyfoot[C]{\thepage}
\renewcommand{\headrulewidth}{0pt}

% --- Soporte para caracteres Unicode adicionales (cuadritos checklist) ---
% Define los simbolos como texto via newunicodechar
\usepackage{newunicodechar}
\newunicodechar{▢}{\ensuremath{\square}}      % cuadro vacio (checkbox)
\newunicodechar{✓}{\ensuremath{\checkmark}}   % check
\newunicodechar{📷}{[\textbf{IMG}]}            % camara -> [IMG]
\usepackage{amssymb}                            % para \checkmark, \square
